\chapter*{Preface}
\addcontentsline{toc}{chapter}{Preface}

Book 16 built a genuine sheaf: the entropy presheaf $\mathcal{S}$ there satisfied the gluing and uniqueness axioms by construction, because the local condition $\diverg \vField = 0$ defining $\mathcal{S}(U)$ is itself local and therefore automatically compatible across overlaps. The present book extends RSVP's sheaf formalism from thermodynamic fields to semantic and cognitive domains, and the extension is not merely a change of subject matter: a presheaf of intelligibility need not glue at all. Where Book 16's local conditions guaranteed coherence, the interpretive structures this book assigns to open regions carry no such guarantee, and it is precisely this difference, sheaves that must cohere against presheaves that may not, that gives this book its subject. Presheaves capture partial understanding and the contextual drift of interpretation, and where Book 16 asked when local thermodynamic data coheres into a global whole, this book asks what happens, and what it means, when local interpretive data does not.

\chapter*{Diagrammatic Overview}
\addcontentsline{toc}{chapter}{Diagrammatic Overview}

The relationship between Book 16's entropy sheaves and this book's semantic presheaves is one of guarantee versus its absence: both are built over the same base manifold $M$ using the same open-set, restriction-map apparatus, but Book 16's sheaf axioms held automatically from the locality of its defining condition, while this book's presheaf of intelligibility $\mathcal{I}$ carries no such automatic guarantee, and Part III below develops exactly what recovering a guarantee, via sheafification, costs.

\chapter*{Reading Note}
\addcontentsline{toc}{chapter}{Reading Note}

Meaning lives in overlaps, and coherence is never global from the start. The reader should not expect this book's presheaves to glue by default, the way Book 16's entropy sheaf did; expect instead to track, chapter by chapter, exactly how much gluing a given interpretive structure achieves and exactly what is lost wherever it falls short.

\part{Defining the Presheaf of Meaning}

\chapter{The Semantic Base Space}

The base manifold $M$ of this book is the identical thermodynamic manifold of experience Book 16 already constructed, with each open set $U \subset M$ now read as representing a cognitive or observational context rather than a purely thermodynamic patch, and morphisms of the base category given, as in that earlier book, by ordinary inclusion $V \hookrightarrow U$. This chapter confirms directly what Book 9's Chapter 10 already assumed when it borrowed this book's vocabulary ahead of its own introduction: the semantic base space is not a separate manifold from Book 16's thermodynamic one but the same $M$, read under a different interpretation of what its open sets represent.

Examples developed here treat overlapping sensory fields and social subdomains as instances of this semantic reading of $M$'s open sets, with the overlaps between such contexts being exactly where Part III's questions about gluing become substantive.

